\section{Introduction}
\subsection{How to read this guide}

This guide explains the whole GNN-AoI project from zero. Read it in order the first time; afterwards use the
contents page to jump to what you need.

\begin{itemize}
  \item \textbf{Section 2} is the glossary: every abbreviation and technical term used anywhere in the project.
  \item \textbf{Section 3} teaches the background ideas (sensor networks, Age of Information, correlation, Kalman
        filters, graph neural networks, reinforcement learning, PPO) from first principles.
  \item \textbf{Section 4} shows how the problem statement came out of the gaps in published research, and lists the
        objectives with how each one was achieved.
  \item \textbf{Section 5} is the complete methodology, equation by equation.
  \item \textbf{Section 6} is the implementation story, step by step, in the order we actually built it.
  \item \textbf{Section 7} explains every design decision: why this dataset, tool or algorithm, and why not the
        alternatives.
  \item \textbf{Section 8} covers the expected output, the obtained output, and every graph and table, with a
        ready-to-use explanation for your faculty.
  \item \textbf{Section 9} gives the conclusion, limitations, future work, likely viva questions and references.
\end{itemize}

Throughout the guide, three kinds of boxes appear:

\begin{whybox}[Why? boxes]
explain the reason behind a choice.
\end{whybox}
\begin{saybox}[How to explain it to your faculty]
give you sentences you can say almost word for word.
\end{saybox}
\begin{notebox}[Note boxes]
point out a subtle detail or an honest limitation.
\end{notebox}

\subsection{The project in one page}

\begin{keybox}[The problem in two sentences]
A base station (the \emph{sink}) must keep track of the readings of many sensors, but the wireless channel lets
only a few of them transmit in each time slot. Which sensors should transmit so that the sink knows the
\emph{whole} field as accurately as possible, while no sensor is ignored for too long?
\end{keybox}

\textbf{The usual answer and why it is not enough.} The standard approach is to poll whichever sensor has gone
longest without reporting. That keeps the \emph{Age of Information} (AoI) low, but it treats every sensor as
independent. In real networks, neighbouring sensors measure almost the same thing (two thermometers in the same room,
two traffic detectors on the same freeway). Polling both wastes a transmission: once one of them reports, the sink
can already infer the other quite well.

\textbf{Our answer.} GNN-AoI is a scheduler that:
\begin{enumerate}
  \item \textbf{learns from historical data} which sensors are correlated (a \emph{correlation graph});
  \item keeps a \textbf{Kalman filter} at the sink, which turns every received reading into an updated estimate of
        \emph{all} sensors, together with how uncertain it is about each one;
  \item uses a \textbf{graph neural network} (GNN) that looks at the graph and the sink's uncertainty and chooses
        which sensors transmit in each slot;
  \item trains that GNN with \textbf{reinforcement learning} (the PPO algorithm), rewarding it for low estimation
        error at the sink;
  \item enforces a \textbf{hard freshness guarantee}: every sensor must report at least once every $F$ slots, no
        matter what the scheduler would prefer.
\end{enumerate}

\textbf{Where we tested it.} On two real sensor datasets, using only time periods that were never used for
training: 45 temperature sensors in the Intel Berkeley Research Lab (2004), and 325 freeway speed sensors in the
San Francisco Bay Area (PEMS-BAY, 2017). We also tested on thousands of simulated sensor layouts.

\textbf{Headline results.}
\begin{center}
\small
\begin{tabular}{@{}lccc@{}}
  \toprule
  & \makecell{\textbf{Intel Lab}\\(45 sensors, $C=6$)} & \makecell{\textbf{PEMS-BAY}\\(325 sensors, $C=40$)} & \makecell{\textbf{Synthetic}\\($M=50$, $C=10$)} \\
  \midrule
  Max-AoI (classic) & 0.292\,$^\circ$C & 2.85 mph & 0.0399 (MSE) \\
  Kalman-greedy & 0.277\,$^\circ$C & 2.95 mph & 0.0205 \\
  \textbf{GNN-PPO (ours)} & \textbf{0.235\,$^\circ$C} & \textbf{2.78 mph} & \textbf{0.0200} \\
  Change vs.\ Max-AoI & $-20\%$ & $-2.4\%$ & $-50\%$ \\
  GNN beats Kalman-greedy in & 10/10 test windows & 15/15 test windows & 73\% of layouts \\
  \bottomrule
\end{tabular}
\end{center}
(Real-data numbers are the sink's root-mean-square error in physical units; lower is better.)

\textbf{What we built.} A tested, open-source Python package (\texttt{gnn\_aoi}) with the simulator, real-data
loaders, Kalman sink, five baseline schedulers, the GNN, the PPO trainer, all experiments and figures, and a live
web demo that replays a real day and shows every transmission as it happens.

\begin{saybox}
``Classic schedulers only ask \emph{who has been silent the longest}. We ask \emph{who would teach the sink the most
right now}, using correlations learned from real data, and we still guarantee that nobody stays silent too long.
On the real Intel lab data this lowers the sink's temperature error by 20 percent with the same channel and the same
freshness guarantee.''
\end{saybox}
